\documentclass[../BTL_PTIT.tex]{subfiles}
\begin{document}

\section{Kết luận và hướng phát triển}

\subsection{Kết luận}
Sau quá trình nghiên cứu và triển khai, nhóm đã xây dựng thành công mô hình hệ thống giám sát chất lượng không khí đáp ứng đầy đủ các mục tiêu đề ra. Hệ thống không chỉ dừng lại ở mức mô hình thử nghiệm mà đã thể hiện được tư duy thiết kế của một sản phẩm IoT hoàn chỉnh.

Các kết quả nổi bật:
\begin{itemize}[label=\textbullet]
	\item Dữ liệu môi trường được cập nhật liên tục với độ trễ thấp qua giao thức MQTT.
	\item Cơ chế điều khiển tự động phản hồi chính xác khi môi trường thay đổi.
	\item Tích hợp thành công các tính năng hiện đại như Cloud Dashboard và cập nhật OTA.
	\item Hệ thống duy trì hoạt động ổn định nhờ cơ chế nguồn dự phòng thông minh.
\end{itemize}

Dù đã đạt được các yêu cầu về chức năng và phi chức năng, hệ thống vẫn còn một số hạn chế về sai số của cảm biến bán dẫn và cần tăng cường thêm các lớp bảo mật dữ liệu trên đường truyền.

\subsection{Hướng phát triển}
Để nâng tầm hệ thống lên quy mô ứng dụng thực tế chuyên sâu hơn, nhóm đề xuất các hướng phát triển:
\begin{itemize}[label=\textbullet]
	\item Nâng cấp độ chính xác: Thay thế cảm biến MQ bằng các dòng cảm biến kỹ thuật số (như NDIR hoặc Laser) để có số liệu chính xác tuyệt đối về bụi mịn và các loại khí đặc thù.
	\item Tối ưu hóa phần mềm: Áp dụng hệ điều hành thời gian thực (FreeRTOS) để quản lý đa tác vụ chuyên nghiệp hơn, bảo đảm tính thời gian thực (real-time) khắt khe trong môi trường công nghiệp.
	\item Trí tuệ nhân tạo (AI): Triển khai các mô hình Machine Learning trên Cloud để dự báo xu hướng ô nhiễm dựa trên dữ liệu lịch sử, thay vì chỉ cảnh báo theo ngưỡng tĩnh như hiện tại.
	\item Bảo mật và quy mô: Tích hợp mã hóa TLS/SSL cho giao thức MQTT và mở rộng hệ thống thành mạng lưới Mesh Network gồm nhiều node giám sát phân tán trong toàn bộ khuôn viên nhà máy.
\end{itemize}

\end{document}